\clearpage
\chapter*{Erratum à l'exposé n\textsuperscript{o} 212}
\addcontentsline{toc}{chapter}{Erratum à l'exposé n\textsuperscript{o} 212}
\setcounter{footnote}{2}
\src{301}

\noindent\textbf{5. TDTE III : Préschémas quotients.}

\begin{center}
[Séminaire Bourbaki, t. 13, 1960/61, n\textsuperscript{o} 212, 20 p.]
\end{center}

\noindent\textbf{Page 15, assertion (i).} --- Le fait que $Y=X/R$ soit
quasi-projectif sur $S$ ne semble prouvé, pour l'instant, que dans le cas où
$R$ provient d'une relation d'équivalence.

\noindent\textbf{Pages 18 et 19.} --- Comme nous le signalons à la fin de
l'exposé suivant, la conjecture envisagée est décidément fausse. Le « fait
positif » signalé dans la remarque 8.1 semble avoir été démontré simultanément
par divers auteurs (NAGATA, ROSENLICHT, GROTHENDIECK, \ldots).

Signalons que l'application faite dans TDTE V de la théorie développée ici du
passage au quotient à la construction des schémas de Picard peut se remplacer
également par une utilisation convenable des schémas de Hilbert (cf. SHMT).
Comme il a été dit dans le § 8, la lacune la plus importante de la théorie,
présentée
\src{302}
ici, est le manque d'un critère d'existence de quotients par une relation
d'équivalence non propre, telles les relations d'équivalence provenant de
certaines opérations du groupe projectif. Un théorème important dans cette
voie a été obtenu par MUMFORD\footnote{MUMFORD (David). --- An elementary
theorem in geometric invariant theory, \emph{Bull. Amer. Math. Soc.}, t. 67,
1961, p. 483--486.}. Pour un raffinement de son résultat, et diverses
applications à la théorie, voir SHMT.
